% -----------------------------------------------------------
\section{Supplement to Concepts}
% -----------------------------------------------------------
	\label{SUPPLEMENT}
	
% -----------------------------------------------------------
% -----------------------------------------------------------
% *A --------------------------------------------------------
% ----------------------------------------------------------- 
% -----------------------------------------------------------

	% ----------------------
	\subsection{Abase} % *ABASE
	% ----------------------
		\label{ABASE}
		
		% -----
		\subsubsection{1828 American Dictionary of the English Language}
		% -----
			
			\begin{compactitem}
				\item[1] The literal sense of \textit{abase} is to lower or depress, to throw or cast down, as used by Bacon, `to \textit{abase} the eye.' But the word is seldom used in reference to material things.
				\item[2] To cast down; to reduce low; to depress; to humble; to degrade; applied to the passions, rank, office, and condition in life.
			\end{compactitem}

	% ----------------------
	\subsection{Abide} % *ABIDE
	% ----------------------
		\label{ABIDE}
		
		% -----
		\subsubsection{1828 American Dictionary of the English Language}
		% -----
		
			% --
			\paragraph{Intransitive verb}
			% --
			
				\begin{compactitem}
					\item[1] To rest, or dwell.
					\item[2] To tarry or stay for a short time.
					\item[3] To continue permanently or in the same state; to be firm and immovable.
					\item[4] To remain, to continue.
				\end{compactitem}
				
			% --
			\paragraph{Transitive verb}
			% --
			
				\begin{compactitem}
					\item[1] To wait for; to be prepared for; to await.
					\item[2] To endure or sustain.
					\item[3] To bear or endure; to bear patiently.
				\end{compactitem}
				
	% ----------------------
	\subsection{Abound} % *ABOUND
	% ----------------------
		\label{ABOUND}
		
		% -----
		\subsubsection{1828 American Dictionary of the English Language}
		% -----
		
			\begin{compactitem}
				\item[1] To have or possess in great quantity; to be copiously supplied; followed by with or in; as to \textit{abound} with provisions; to \textit{abound} in good things.
				\item[2] To be in great plenty; to be very prevalent.
			\end{compactitem}
			
	% ----------------------
	\subsection{Abroad} % *ABROAD
	% ----------------------
		\label{ABROAD}
		
		% -----
		\subsubsection{1828 American Dictionary of the English Language}
		% -----
		
			\begin{compactitem}
				\item In a general sense, at large; widely; not confined to narrow limits. Hence,
				\item[1] In the open air.
				\item[2] Beyond or out of the walls of a house, as to walk abroad.
				\item[3] Beyond the limits of a camp.
				\item[4] Beyond the bounds of a country; in foreign countries - as to go abroad for an education.
				\item[5] Extensively; before the public at large.
				\item[6] Widely; with expansion; as a tree spreads its branches abroad
			\end{compactitem}
			
	% ----------------------
	\subsection{Adam-ondi-Ahman} % *ADAM-ONDI-AHMAN *ADAMONDIAHMAN
	% ----------------------
		\label{ADAM-ONDI-AHMAN}
		
		% -----
		\subsubsection{Scriptural Usage}
		% -----
		
			\begin{tabular}{ll}
				Total occurrences of ``Adam-ondi-Ahman'' in the scriptures & 5 \\
				Total occurrences in the DC & 5
			\end{tabular}
				
			\begin{compactdesc}
				\item[{\hyperref[DC78-15]{DC 78:15}}]
				\item[{\hyperref[DC107-53]{DC 107:53}}]
				\item[{\hyperref[DC116-1]{DC 116:1}}]
				\item[{\hyperref[DC117-8]{DC 117-8}}]
				\item[{\hyperref[DC117-11]{DC 117-11}}]
			\end{compactdesc}
			
	% ----------------------
	\subsection{Although} % *ALTHOUGH
	% ----------------------
		\label{ALTHOUGH}
		
		% -----
		\subsubsection{1828 American Dictionary of the English Language}
		% -----
		
			\begin{compactitem}
				\item Grant all this; be it so; allow all; suppose that; admit all that; as, `although the fig tree shall not blossom.' Habakkuk 3:17. That is, grant, admit or suppose what follows---`the fig tree shall not blossom.' It is a transitive verb, and admits after it the definitive that---although that the fig tree shall not blossom; but this use of the verb, has been long obsolete. The word may be defined by notwithstanding, non obstante; as not opposing may be equivalent to admitting or supposing.
			\end{compactitem}
			
	% ----------------------
	\subsection{Appendange} % *APPENDAGE
	% ----------------------
		\label{APPENDAGE}
		
		% -----
		\subsubsection{1828 American Dictionary of the English Language}
		% -----
		
			\begin{compactitem}
				\item Something added to a principal or greater thing, though not necessary to it, as a portico to a house.
			\end{compactitem}
		
		% -----
		\subsubsection{Oxford English Dictionary}
		% -----
		
			\begin{compactitem}
				\item That which is attached as if by being hung on; a subsidiary external adjunct, addition, or accompaniment, which does not form an essential part of that to which it is added, but is usually natural or appropriate to it.
				\item[1] of things material.
				\item[1.a] Generally.
				\item[1.b] \textit{esp.} An addition to territory or property.
				\item[1.c] An addition in writing; an appendix.
				\item[1.d] \textit{Natural History.} A subordinate or subsidiary organ.
				\item[2] of things immaterial.
				\item[3] \textit{transferred} of persons.
			\end{compactitem} 

% -----------------------------------------------------------
% -----------------------------------------------------------
% *B --------------------------------------------------------
% ----------------------------------------------------------- 
% -----------------------------------------------------------

	% ----------------------
	\subsection{Betimes} % *BETIMES
	% ----------------------
		\label{BETIMES}
		
		% -----
		\subsubsection{1828 American Dictionary of the English Language}
		% -----
		
			\begin{compactitem}
				\item[1] Seasonably; in good season or time; before it is late.
				\item[2] Soon; in a short time.
			\end{compactitem}
		
		% -----
		\subsubsection{Oxford English Dictionary}
		% -----
		
			\begin{compactitem}
				\item[1] At an early time, period, or season; early in the year; early in life.
				\item[2] \textit{spec.} At an early hour, early in the morning.
				\item[3] In good time, in due time; while there is yet time, before it is too late.
				\item[4] In a short time, soon, speedily, anon, forthwith.
			\end{compactitem}
			
	% ----------------------
	\subsection{Bishopric} % *BISHOPRIC
	% ----------------------
		\label{BISHOPRIC}
		
		% -----
		\subsubsection{1828 American Dictionary of the English Language}
		% -----
		
			\begin{compactitem}
				\item[1] A diocese; the district over which the jurisdiction of bishop extends. in England, are twenty-four bishoprics, besides that of Sodor and Man; in Ireland, eighteen.
				\item[2] The charge of instructing and governing in spiritual concerns; office. \hyperref[ACTS-1-20]{Acts 1:20}.
			\end{compactitem}
			
	% ----------------------
	\subsection{Book (as a heavenly records)} % *BOOK
	% ----------------------
		\label{BOOK}
		
		Here's a list of references to such books:
		
		\begin{compactitem}
			\item ``the book of the names of the sanctified, even them of the celestial world'' (\hyperref[DC88-2]{DC 88:2})
			\item ``the volume of the book,'' in which deeds committed against God are written (\hyperref[DC99-5]{DC 99:5})
			\item \hyperref[DC128-6]{DC 128:6--7}, referring to \hyperref[REVELATIONNT-20-12]{Revelation 20:12}
		\end{compactitem}

% -----------------------------------------------------------
% -----------------------------------------------------------
% *C --------------------------------------------------------
% ----------------------------------------------------------- 
% -----------------------------------------------------------

	% ----------------------
	\subsection{Caleb} % *CALEB
	% ----------------------
		\label{CALEB}
		
		% -----
		\subsubsection{NIEBC, ``Caleb'' \#1}
		% -----
		
			``Caleb, from the tribe of Judah, was the son of Jephunneh the Kenizzite. Moses asked him to represent his tribe in spying out the land of Canaan (Num. 13:6). This he did with twelve others. The land they saw was very fertile. They came back with large bunches of grapes and great reports of the land's wealth and prosperity. However, the land was found to be occupied by some fearful people, especially the descendants of the giant Anak, who lived around Hebron. Ten of the spies came back so disillusioned and so convinced they could never win in battle against these people that they turned on Moses. The Israelites, on hearing this, argued that they would have been better off back in Egypt (Num. 13:26--29; 13:31--14:4).
			
			``Of all the spies, only Caleb and Joshua had sufficient faith in God to know that he could enable them to take the land of Canaan. First they gave an alternative report to Moses, and then they appealed to the people. Their complete trust in the sovereign power of God is proclaimed aloud to the people in their speech: `If the LORD is pleased with us, he will lead us into that land, a land flowing with milk and honey, and will give it to us. Only do not rebel against the LORD. And do not be afraid of the people . . . Their protection is gone, but the LORD is with us' (Num. 14:7--9).
			
			``As a judgment on the people for their lack of faith, that whole generation was never allowed to enter Canaan, but rather spent the time `wandering' in the wilderness (Num. 14:22--23; 26:65). Only Caleb and Joshua were allowed to live long enough to enter the land, and both had a very significant part in the battles to take Canaan many years later (Num. 14:24, 30). It is of particular significance that Caleb, who so many years before had believed that the LORD would destroy the Anakites, was the leader of the attack on the Hebron area. He personally led the forces that defeated the giants, and as a result he was later given that area as an inheritance by Joshua (Josh. 14:6--15; 15:13--14). Caleb offered his daughter Acsah in marriage to the man who would conquer Kiriath-sepher (Debir) for him. Othniel, Caleb's nephew, did this and so married her (Josh. 15:16--18). His sons are listed in 1 Chron. 4:15.
			
			``Caleb's 'wholehearted faith' (Josh. 14:14; Deut 1:36 etc.) became an example for future generations. This faith issued in a practical and calm reliance upon the LORD God as the one who was all-powerful and whose will could not be thwarted. To believe this would be one thing, but to act on it in the face of such adversaries was of the essence of the true faith which Caleb demonstrated and from which all can learn.''
			
		% -----
		\subsubsection{NIEBC, ``Caleb'' \#2}
		% -----
		
			``This Caleb was the brother of Jerahmeel and the son of Hezron. He was of the tribe of Judah and married to Azubah. When she died he married Ephrath, by whom he had a son called Hur. His other sons were Mesha and Mareshah (1 Chron. 2:9, 18--19, 42, 50). Hur's grandson Bezalel was a craftsman who worked on the Tabernacle.''
			
	% ----------------------
	\subsection{Carmel, Mount} % *CARMEL *MOUNTCARMEL
	% ----------------------
		\label{CARMEL}
		\label{MOUNTCARMEL}
		
		% -----
		\subsubsection{BIBD, ``Carmel, Mount''}
		% -----
		
			``The wooded mountain promontory on the Mediterranean, near modern Haifa. The name means `the garden.' It forms a northern barrier to the coastal plain of Sharon. Mount Carmel provided the perfect stage for its most significant event, the confrontation between Elijah and the prophets of Baal (1 Kings 18), the god of storms and therefore agricultural produce. The mountain's high elevation meant that it was lush until a drought. When the prophets threatened that Carmel would wither, conditions were extreme (Isa. 33:9; Amos 1:2; Nah.1:4).''
			
		% -----
		\subsubsection{ZIBD, ``Carmel''}
		% -----
		
			``When the word occurs with the definite article, it generally refers to Mount Carmel and is often used to illustrate a beautiful and fruitful place (Isa. 35:2; but see 33:9, which pictures God's judgment). South of Carmel lies the fruitful plain of SHARON and NE of it flows the river KISHON through the Plain of ESDRAELON. At Carmel, ELIJAH stood against 450 heathen prophets and defeated them (1 Ki. 18). ELISHA also visited Carmel (2 Ki. 2:25; 4:25).''
			
	% ----------------------
	\subsection{Cavity} % *CAVITY
	% ----------------------
		\label{CAVITY}
		
		% -----
		\subsubsection{1828 American Dictionary of the English Language}
		% -----

			\begin{compactitem}
				\item A hollow place; hollowness; an opening; as the \textit{cavity} of the mouth or throat. This is a word of very general signification.
			\end{compactitem}
			
	% ----------------------
	\subsection{Circumscribe} % *CIRCUMSCRIBE
	% ----------------------
		\label{CIRCUMSCRIBE}
		
		
		% -----
		\subsubsection{1828 American Dictionary of the English Language}
		% -----
		
			\begin{compactitem}
				\item[1] To inclose within a certain limit; to limit, bound, confine.
				\item[2] To write round.
			\end{compactitem}
			
		% -----
		\subsubsection{Oxford English Dictionary}
		% -----
		
			\begin{compactitem}
				\item[1.a] \textit{transitive.} To draw a line round; to encompass with (or as with) a bounding line, to form the boundary of, to bound.
				\item[1.b] To encompass (without a line), to encircle.
				\item[2.a] To mark out or lay down the limits of; to enclose within limits, limit, bound, confine (usually \textit{figurative}); \textit{esp.} to confine within narrow limits, to restrict the free or extended action of, to hem in, restrain, abridge.
				\item[2.b] To mark off, to define logically.
				\item[3.a] \textit{Geometry.} To describe (a figure) about another figure so as to touch it at certain points or parts without cutting it.
				\item[3.b] With the figure so described as subject of the verb.
				\item[4.a] To write or inscribe around (a coin, etc., with an inscription, or an inscription on or about a coin, etc.). \textit{Obsolete.}
				\item[4.b] To join in signing a `round-robin'.
			\end{compactitem}
			
	% ----------------------
	\subsection{Cleave} % *CLEAVE
	% ----------------------
		\label{CLEAVE}
		
		% -----
		\subsubsection{1828 American Dictionary of the English Language}
		% -----
		
			% --
			\paragraph{Intransitive verb}
			% --
		
				\begin{compactitem}
					\item[1] To stick; to adhere; to hold to.
					\item[2] To unite aptly; to fit; to sit well on.
					\item[3] To unite or be united closely in interest or affection; to adhere with strong attachment.
				\end{compactitem}
				
			% --
			\paragraph{Transitive verb}
			% --
			
				\begin{compactitem}
					\item[1] To part or divide by force; to split or rive; to open or serve the cohering parts of a body, by cutting or by the application of force; as, to \textit{cleave} wood; to \textit{cleave} a rock; to \textit{cleave} the flood.
					\item[2] To part or open naturally.
				\end{compactitem}
				
			% --
			\paragraph{Intransitive verb}
			% --
			
				\begin{compactitem}
					\item To part; to open; to crack; to separate, as parts of cohering bodies; as, the ground cleaves by frost.
				\end{compactitem}
			
	% ----------------------
	\subsection{Constrain} % *CONSTRAIN
	% ----------------------
		\label{CONSTRAIN}
		
		% -----
		\subsubsection{1828 American Dictionary of the English Language}
		% -----
		
			\begin{compactitem}
				\item In a general sense, to strain; to press; to urge; to drive; to exert force, physical or moral, either in urging to action or in restraining it. Hence,
				\item[1] To compel or force; to urge with irresistible power, or with a power sufficient to produce the effect.
				\item[2] To confine by fore; to restrain from escape or action; to repress.
				\item[3] To hold by force; to press; to confine.
				\item[4] To constringe; to bind.
				\item[5] To tie fast; to bind; to chain; to confine.
				\item[6] To necessitate.
				\item[7] To force; to ravish. [\textit{Not used.}]
				\item[8] To produce in opposition to nature; as a constrained voice; constrained notes.
			\end{compactitem}
			
	% ----------------------
	\subsection{Council} % *COUNCIL
	% ----------------------
		\label{COUNCIL}
		
		% -----
		\subsubsection{1828 American Dictionary of the English Language}
		% -----
		
			% --
			\paragraph{Noun}
			% --
			
				(Note that there is no 1828 entry for a verb spelled ``council.'' The beginning of ``council'' as a noun says, ``This word is often confounded with `counsel,' with which it has no connection. `Council' is a collection or assembly.'' The OED also does not have a verb entry for ``council.'')
				
				\begin{compactitem}
					\item An assembly of men summoned or convened for consultation, deliberation and advice.
					\item A body of men specially designated to advise a chief magistrate in the administration of the government, as in Great Britain.
					\item In some of the American states, a branch of the legislature, corresponding with the senate in other states, and called legislative \textit{council}.
					\item An assembly of prelates and doctors, convened for regulating matters of doctrine an discipline in the church.
					\item Act of deliberation; consultation of a \textit{council}.
				\end{compactitem}
			
	% ----------------------
	\subsection{Counsel} % *COUNSEL
	% ----------------------
		\label{COUNSEL}
		
		% -----
		\subsubsection{1828 American Dictionary of the English Language}
		% -----
		
			% --
			\paragraph{Transitive verb}
			% --
		
			\begin{compactitem}
				\item To give advice or deliberate opinion to another for the government of his conduct; to advise.
				\item To exhort, warn, admonish, or instruct. We ought frequently to \textit{counsel} our children against the vices of the age.
				\item To advise or recommend; as, to \textit{counsel} a crime. [Not much used.]
			\end{compactitem}
			
	% ----------------------
	\subsection{Course} % *COURSE
	% ----------------------
		\label{COURSE}
		
		% -----
		\subsubsection{1828 American Dictionary of the English Language}
		% -----
		
			\begin{compactitem}
				\item In its general sense, a passing; a moving, or motion forward, in a direct or curving line; applicable to any body or substance, solid or fluid. Applied to animals, a running, or walking; a race; a career; a passing, or passage, with any degree of swiftness indefinitely. Applied to fluids, a flowing, as in a stream in any direction; as a straight \textit{course} or winding \textit{course} It is applied to water or other liquids, to air or wind, and to light, in the sense of motion or passing. Applied to solid bodies, it signifies motion or passing; as the \textit{course} of a rolling stone; the \textit{course} of a carriage; the \textit{course} of the earth in its orbit. Applied to navigation, it signifies a passing or motion on water, or in balloons in air; a voyage.
				\item The direction of motion; line of advancing; point of compass, in which motion is directed; as, what \textit{course} shall the pilot steer? In technical language, the angel contained between the nearest meridian and that point of compass on which a ship sails in any direction.
				\item Ground on which a race is run.
				\item A passing or process; the progress of any thing; as the \textit{course} of an argument, or of a debate; a \textit{course} of thought or reflexion.
				\item Order of proceeding or of passing from an ancestor to an heir; as the \textit{course} of descent in inheritance.
				\item Order; turn; class; succession of one to another in office, or duty. The chief fathers of every \textit{course} 1 Chronicles 27:1. Solomon appointed the \textit{course}s of the priests. 2 Chronicles 8:14.
				\item Stated and orderly method of proceeding; usual manner. He obtained redress in due \textit{course} of law. Leave nature to her \textit{course}.
				\item Series of successive and methodical procedure; a train of acts, or applications; as a \textit{course} of medicine administered.
				\item A methodical series, applied to the arts or sciences; a systemized order of principles in arts or sciences, for illustration of instruction. We say, the author has completed a \textit{course} of principles or of lectures in philosophy. Also, the order pursued by a student; as, he has completed a \textit{course} of studies in law or physics.
				\item Manner of proceeding; way of life or conduct; deportment; series of actions. That I might finish my \textit{course} with joy. Acts 20:24. Their \textit{course} is evil. Jeremiah 23:10.
				\item Line of conduct; manner of proceeding; as, we know not what \textit{course} to pursue.
				\item Natural bent; propensity; uncontrolled will. Let not a perverse child take his own \textit{course}.
				\item Tilt; act of running in the lists.
				\item Orderly structure; system. The tongue setteth on fire the \textit{course} of nature. James 3:6.
				\item Any regular series. In architecture, a continued range of stones, level or of the same highth, throughout the whole length of the building, and not interrupted by any aperture. A laying of bricks, etc. 
				\item The dishes set on table at one time; service of meat. 
				\item Regularity; order; regular succession; as, let the classes follow in \textit{course} .
				\item Empty form; as, compliments are often words of \textit{course}. Of \textit{course} by consequence; in regular or natural order; in the common manner of proceeding; without specila direction or provision. This effect will follow of \textit{course} If the defendant resides no in the state, the cause is continued of \textit{course}.
			\end{compactitem}

% -----------------------------------------------------------
% -----------------------------------------------------------
% *D --------------------------------------------------------
% ----------------------------------------------------------- 
% -----------------------------------------------------------

	% ----------------------
	\subsection{Dispensation} % *DISPENSATION
	% ----------------------
		\label{DISPENSATION}
		
		% -----
		\subsubsection{1828 American Dictionary of the English Language}
		% -----
		
			\begin{compactitem}
				\item[1] Distribution; the act of dealing out to different persons or places; as the \textit{dispensation} of water indifferently to all parts of the earth.
				\item[2] The dealing of God to his creatures; the distribution of good and evil, natural or moral, in the divine government.
				\item[3] The granting of a license, or the license itself, to do what is forbidden by laws or canons, or to omit something which is commanded; that is, the dispensing with a law or canon, or the exemption of a particular person from the obligation to comply with its injunctions. The pope has power to dispense with the canons of the church, but has no right to grant \textit{dispensation}s to the injury of a third person.
				\item[4] That which is dispensed or bestowed; a system of principles and rites enjoined; as the Mosaic \textit{dispensation}; the gospel \textit{dispensation}; including, the former the Levitical law and rites; the latter the scheme of redemption by Christ.
			\end{compactitem}

	% ----------------------
	\subsection{Distill} % *DISTIL *DISTILL
	% ----------------------
		\label{DISTIL} \label{DISTILL}
		
		``Distil'' (spelled that way) appears three times in the scriptures: \hyperref[DEUTERONOMY-32-2]{Deuteronomy 32:2}; \hyperref[JOB-36-28]{Job 36:28}; and \hyperref[DC121-45]{DC 121:45}.
		
		% -----
		\subsubsection{1828 American Dictionary of the English Language}
		% -----
		
			% --
			\paragraph{Intransitive verb}
			% --
		
				\begin{compactitem}
					\item[1] To drop; to fall in drops.
					\item[2] To flow gently, or in a small stream.
					\item[3] To use a still; to practice distillation.
				\end{compactitem}
				
			% --
			\paragraph{Transitive verb}
			% --
			
				\begin{compactitem}
					\item[1] To let fall in drops; to throw down in drops. The clouds distill water on the earth.
					\item[2] To extract by heat; to separate spirit or essential oils from liquor by heat or evaporation, and convert that vapor into a liquid by condensation in a refrigeratory; to separate the volatile parts of a substance by heat; to rectify; as, to distill brandy from wine, or spirit form melasses.
					\item[3] To extract spirit from, by evaporation and condensation; as, to distill cyder or melasses; to distill wine.
					\item[4] To extract the pure part of a fluid; as, to distill water.
					\item[5] To dissolve or melt. [Unusual.]
				\end{compactitem}
			
		% -----
		\subsubsection{Oxford English Dictionary}
		% -----
			
			\begin{compactitem}
				\item[1.a] \textit{intransitive.} To trickle down or fall in minute drops, as rain, tears; to issue forth in drops or in a fine moisture; to exude.
				\item[1.b] To pass or flow gently. Chiefly \textit{figurative.}
				\item[1.c] To melt into, or become dissolved in, \textit{tears}.
				\item[1.d] To drip or be wet \textit{with}.
				\item[2] \textit{transitive.} To let fall or give forth in minute drops, or in a vapour which condenses into drops.
				\item[3] \textit{transferred} and \textit{figurative.} To give forth or impart in minute quantities; to infuse; to instil.
				\item[4.a] To subject to the process of distillation; to vaporize a substance by means of heat, and then condense the vapour by exposing it to cold, so as to obtain the substance or one of its constituents in a state of concentration or purity. Primarily said of a liquid, the vapour of which when condensed is again deposited in minute drops of pure liquid; but extended also to the volatilizing of solids, the products of which may be gaseous.
				\item[4.b] To extract the essence of (a plant, etc.) by distillation; to obtain an extract of.
				\item[4.c] To transform or convert (into something) by distillation. Also \textit{figurative.}
				\item[4.d] \textit{absol.} To perform distillation.
				\item[4.e] \textit{figurative.} To extract the quintessence of; to concentrate, purify.
				\item[4.f] To drive (a volatile constituent) \textit{off} or \textit{out} by distillation. \textit{Also figurative.}
				\item[5.a] To obtain, extract, produce, or make, by distillation.
				\item[5.b] \textit{figurative.}
				\item[6] To become vaporized and then condensed into liquid; to undergo distillation; to drop, pass, or condense from the still. \textit{\textbf{to distil over}}: to pass over in the form of vapour which again condenses into a liquid.
				\item[7] \textit{transitive.} To melt, dissolve (\textit{literal} and \textit{figurative}). \textit{Obsolete.}
			\end{compactitem}
			
	% ----------------------
	\subsection{Draw} % *DRAW
	% ----------------------
		\label{DRAW}
		
		% -----
		\subsubsection{Greek New Testament}
		% -----
		
			%
			\subparagraph{\texorpdfstring{\gr{ἕλκω}}{}}
			%
				\label{HELKO}
				
				% \begin{tabular}{ll}
					% Total occurrences in the NT & 
				% \end{tabular}
				
				\begin{compactdesc}
					\item[Some NT uses] \textbf{} % Change to "all" if they're all listed
						\begin{compactdesc}
							\item[{\hyperref[JOHN-6-44]{John 6:44}}]
							\item[{\hyperref[JOHN-12-32]{John 12:32}}]
						\end{compactdesc}
				\end{compactdesc}
			
				\begin{compactdesc}
				
					\item[BDAG 281a] \phantom{.}
						\begin{compactitem}
							\item[1] \textbf{to move an object from one area to another in a pulling motion, \textit{draw}}, with implication that the object being moved is incapable of propelling itself or in the case of person is unwilling to do so voluntarily, in either case with implication of exertion on the part of the mover...
								\begin{compactitem}
									\item See meanin \#2 below; the sense of being unwilling is clear in the Plato and Epictetus quotations below.
								\end{compactitem}
							\item[2] \textbf{to draw a person in the direction of values for inner life, \textit{draw, attract}}, an extended figurative use of meaning \#1
								\begin{compactitem}
									\item John 6:44 is listed here.
									\item The meaning is not necessarily positive---the entry references Plato's \textit{Phaedrus} 237e--238a: ``Now when judgment is in control and leads us by reasoning toward what is best, that sort of self-control is called `being in your right mind'; but when desire takes command in us and drags [\gr{ἕλκω}] us without reasoning toward pleasure, then its command is known as `outrageousness'.'' It looks like the possible unwillingness involved in meaning \#1 also applies here.
									\item The entry also references Epictetus, \textit{Discourses} 2.20.15: ``So what was it that roused Epicurus from his sleep and made him write what he wrote? It must have been what is, after all, the strongest force in human life, nature, which compels [\gr{ἕλκω}] people to do its bidding, reluctant and complaining though they may be.''
									\item It reminds me of the fact that the Holy Ghost ``entices'' us, or tries to convince or ``draw'' us, to the good. See \hyperref[MOSIAH-3-19]{Mosiah 3:19}.
								\end{compactitem}
							\item[3] \textbf{to appear to be pulled in a certain direction, \textit{flow}} an extended figurative use
						\end{compactitem}
						
					% \item[Brill , PDF ] \phantom{.}
						% \begin{compactitem}
							% \item 
						% \end{compactitem}
						
					% \item[CGL , PDF ] \phantom{.}
						% \begin{compactitem}
							% \item 
						% \end{compactitem}
						
					% \item[LSJ , PDF ] \phantom{.}
						% \begin{compactitem}
							% \item 
						% \end{compactitem}
						
					% \item[TDNT ] \phantom{.}
						% \begin{compactitem}
							% \item 
						% \end{compactitem}
						
					% \item[TDNTA ] \phantom{.}
						% \begin{compactitem}
							% \item 
						% \end{compactitem}
						
					% \item[NIDNTT ] \phantom{.}
						% \begin{compactitem}
							% \item 
						% \end{compactitem}
						
					% \item[NIDNTTE ] \phantom{.}
						% \begin{compactitem}
							% \item 
						% \end{compactitem}
						
				\end{compactdesc}

% -----------------------------------------------------------
% -----------------------------------------------------------
% *E --------------------------------------------------------
% ----------------------------------------------------------- 
% -----------------------------------------------------------

	% ----------------------
	\subsection{Effectual} % *EFFECTUAL
	% ----------------------
		\label{EFFECTUAL}
		
		``Effectual'' appears four times in the scriptures---once in the NT (\hyperref[1CORINTHIANS-16-9]{1 Corinthians 16:9}), and three times in the DC \hyperref[DC100-3]{100:3}, \hyperref[DC112-19]{112:19}, \hyperref[DC118-3]{118:3}).
		
		% -----
		\subsubsection{1828 American Dictionary of the English Language}
		% -----
		
			\begin{compactitem}
				\item Producing an effect, or the effect desired or intended; or having adequate power or force to produce the effect.
				\item[1] Veracious; expressive of facts. [Not used.]
				\item[2] \textit{effectual} assassin, in Mitford, is unusual and not well authorized.
			\end{compactitem}
			
	% ----------------------
	\subsection{Essay} % *ESSAY
	% ----------------------
		\label{ESSAY}
		
		% -----
		\subsubsection{1828 American Dictionary of the English Language}
		% -----
		
			\begin{compactitem}
				\item[1] To try; to attempt; to endeavor; to exert one's power or faculties, or to make an effort to perform any thing.
				\item[2] To make experiment of.
				\item[3] To try the value and purity of metals. In this application, the word is now more generally written assay, which see.
			\end{compactitem}

	% ----------------------
	\subsection{Establish} % *ESTABLISH
	% ----------------------
		\label{ESTABLISH}
		
		% -----
		\subsubsection{1828 American Dictionary of the English Language}
		% -----
		
			\begin{compactitem}
				\item[1] To set and fix firmly or unalterably; to settle permanently.
				\item[2] To found permanently; to erect and fix or settle; as, to establish a colony or an empire.
				\item[3] To enact or decree by authority and for permanence; to ordain; to appoint; as, to establish laws, regulations, institutions, rules, ordinances, etc.
				\item[4] To settle or fix; to confirm; as, to establish a person, society or corporation, in possessions or privileges.
				\item[5] To make firm; to confirm; to ratify what has been previously set or made.
				\item[6] To settle or fix what is wavering, doubtful or weak; to confirm.
				\item[7] To confirm; to fulfill; to make good.
				\item[8] To set up in the place of another and confirm.
			\end{compactitem}
			
	% ----------------------
	\subsection{Exactness} % *EXACTNESS
	% ----------------------
		\label{EXACTNESS}
		
		% -----
		\subsubsection{1828 American Dictionary of the English Language}
		% -----
		
			\begin{compactitem}
				\item Accuracy; nicety; precision; as, to make experiments with \textit{exactness}.
				\item[1] Regularity; careful conformity to law or rules of propriety; as \textit{exactness} of deportment.
				\item[2] Careful observance of method and conformity to truth; as \textit{exactness} in accounts or business.
			\end{compactitem}

	% ----------------------
	\subsection{Expedient} % *EXPEDIENT
	% ----------------------
		\label{EXPEDIENT}
		
		% -----
		\subsubsection{1828 American Dictionary of the English Language}
		% -----
		
			\begin{compactitem}
				\item[1] Literally, hastening; urging forward. Hence, tending to promote the object proposed; fit or suitable for the purpose; proper under the circumstances. Many things may be lawful, which are not \textit{expedient}
				\item[2] Useful; profitable.
				\item[3] Quick; expeditious. [Not used.]
			\end{compactitem}
			
		% -----
		\subsubsection{Oxford English Dictionary}
		% -----

			\begin{compactitem}
				\item[I] Senses relating to speed or agility.
				\item[I.1.a] Hasty, `expeditious', speedy. Also, of a march: Direct. Obsolete.
				\item[I.1.b] Quasi-adv. Nimbly, skilfully. Obsolete.
				\item[II] Senses relating to advantage or suitability.
				\item[II.2] Conducive to advantage in general, or to a definite purpose; fit, proper, or suitable to the circumstances of the case. Const. for, to.
				\item[II.2.a] As predicate or complement, often with subject it, and followed by infinitive phrase or noun-sentence.
				\item[II.2.b] Qualifying a noun.
				\item[II.3] In depreciative sense, `useful' or `politic' as opposed to `just' or `right'. Often absol.
				\item[II.4] Studious of `expediency'.
			\end{compactitem}
			
		% -----
		\subsubsection{Expedient and pragmatic}
		% -----
			
			When you look at the \hyperref[EXPEDIENT-1828]{definitions} of ``expedient'' in the 1828 dictionary, the word definitely comes out looking like a word of pragmatism. It's about doing what is needful under the circumstances, to promote a certain end.
			
	% ----------------------
	\subsection{Express} % *EXPRESS
	% ----------------------
		\label{EXPRESS}
		
		% -----
		\subsubsection{1828 American Dictionary of the English Language}
		% -----
		
			\begin{compactitem}
				\item Plain; clear; expressed; direct not ambiguous.
				\item[1] Given in direct terms; not implied or left to inference. 
				\item[2] Copied; resembling; bearing an exact representation.
				\item[3] Intended or sent for a particular purpose, or on a particular errand; as, to send a messenger \textit{express}.
			\end{compactitem}
			
		% -----
		\subsubsection{Oxford English Dictionary}
		% -----

			\begin{compactitem}
				\item[1.a] Of a meaning, rule, etc.: clearly and unambiguously stated; explicit. Of written or spoken language: formulated in clear and distinct terms; unmistakable in meaning.
				\item[1.b] Expressing oneself in a clear, direct, or unreserved manner. Also, of a state of mind: fixed, determined. \textit{Obsolete.}
				\item[1.c] Of a voice: clear, distinct. \textit{Obsolete. rare.}
				\item[1.d] Specifically designated or considered; special. Now chiefly of a purpose, intention, etc.
				\item[2] (legal use)
				\item[3.a] Specially designed or intended for a particular aim, goal, or object; done, made, or sent on purpose or deliberately.
				\item[3.b] (railway use)
				\item[3.c] Generally: characterized by greater speed or rapidity than usual; (very) fast.
				\item[3.d] (postal use)
				\item[3.e] (transportation use)
				\item[4] Of an image, form, etc.: truly depicting the original, exactly resembling, exact.
			\end{compactitem}

% -----------------------------------------------------------
% -----------------------------------------------------------
% *F --------------------------------------------------------
% ----------------------------------------------------------- 
% -----------------------------------------------------------

	% ----------------------
	\subsection{Forget} % *FORGET
	% ----------------------
		\label{FORGET}
		
		% -----
		\subsubsection{Oxford English Dictionary}
		% -----
		
			\begin{compactitem}
				\item[1.a] To lose remembrance of; to cease to retain in one's memory.
				\item[1.b] To fail to recall to mind; not to recollect.
				\item[1.d] \textit{absol.} (or \textit{intransitive}) Also, to forget about: not to recall the facts concerning; not to remember to take action in the matter of (\textit{colloquial}).
				\item[2.a] To omit or neglect through inadvertence.
				\item[2.b] To omit to take, leave behind inadvertently.
				\item[2.d] To omit to mention, leave unnoticed, pass over inadvertently.
				\item[3.a] To cease or omit to think of, let slip out of the mind, leave out of sight, take no note of.
				\item[4] In stronger sense: To neglect wilfully, take no thought of, disregard, overlook, slight.
			\end{compactitem}

	% ----------------------
	\subsection{Foundation} % *FOUNDATION
	% ----------------------
		\label{FOUNDATION}
		
		% -----
		\subsubsection{Greek New Testament}
		% -----
		
			% --
			\paragraph{\texorpdfstring{\gr{καταβολή}}{}}
			% --
				\label{KATABOLE}
				
				% \begin{tabular}{ll}
					% Total occurrences in the NT & 
				% \end{tabular}
				
				\begin{compactdesc}
					\item[Some NT uses] \textbf{} % Change to "all" if they're all listed
						\begin{compactdesc}
							\item[{\hyperref[MATTHEW-13-35]{Matthew 13:35}}]
						\end{compactdesc}
				\end{compactdesc}
			
				\begin{compactdesc}
				
					\item[BDAG 456b] \textbf{}
						\begin{compactitem}
							\item[1] \textbf{the act of laying something down with implication of providing a base for something, \textit{foundation}}. Readily connected with the idea of founding is the sense \textit{beginning}...\gr{ἀπὸ καταβολῆς κόσμου} \textit{from the foundation of the world}...
								\begin{compactitem}
									\item Note that the SBLG for Matthew 13:35 lacks \gr{κόσμου}, but it's not hard to think that that's the intended meaning.
								\end{compactitem}
							\item[2] \gr{κ.}, a technical term for the \textbf{sowing} of seed, used of begetting...If this meaning is correct for \hyperref[HEBREWS-11-11]{Hb 11:11}, there is probably some error in the text, since this expression could not be used of Sarah, but only of Abraham...
						\end{compactitem}
						
					% \item[Brill , PDF ] \textbf{}
						% \begin{compactitem}
							% \item 
						% \end{compactitem}
						
					% \item[CGL , PDF ] \textbf{}
						% \begin{compactitem}
							% \item 
						% \end{compactitem}
						
					% \item[LSJ , PDF ] \textbf{}
						% \begin{compactitem}
							% \item 
						% \end{compactitem}
						
					% \item[TDNT ] \textbf{}
						% \begin{compactitem}
							% \item 
						% \end{compactitem}
						
					% \item[TDNTA ] \textbf{}
						% \begin{compactitem}
							% \item 
						% \end{compactitem}
						
					% \item[NIDNTT ] \textbf{}
						% \begin{compactitem}
							% \item 
						% \end{compactitem}
						
					% \item[NIDNTTE ] \textbf{}
						% \begin{compactitem}
							% \item 
						% \end{compactitem}
						
				\end{compactdesc}

	% ----------------------
	\subsection{Forsake} % *FORSAKE
	% ----------------------
		\label{FORSAKE}
		
		% -----
		\subsubsection{1828 American Dictionary of the English Language}
		% -----
		
			\begin{compactitem}
				\item[1] To quit or leave entirely; to desert; to abandon; to depart from. Friends and flatterers \textit{forsake} us in adversity.
				\item[2] To abandon; to renounce; to reject.
				\item[3] To leave; to withdraw from; to fail. In anger, the color forsakes the cheeks. In severe trials, let not fortitude \textit{forsake} you.
				\item[4] In scripture, God \textit{forsakes} his people, when he withdraws his aid, or the light of his countenance.
			\end{compactitem} 
			
		% -----
		\subsubsection{Oxford English Dictionary}
		% -----
		
			\begin{compactitem}
				\item[1.a] To deny (an accusation, an alleged fact, etc.). \textit{Obsolete.}
				\item[1.b] To deny knowledge of (a person). \textit{Obsolete.}
				\item[1.c] To deny, renounce, or repudiate allegiance to (God, a lord, etc.). 
				\item[1.d] To `deny' (oneself).
				\item[2.a] To decline or refuse (something offered). With simple object or to and infinitive. \textit{Obsolete.}
				\item[2.b] To decline or refuse to bear, encounter, have to do with, undertake; to avoid, shun. \textit{Obsolete.}
				\item[2.c] To refuse respect or obedience to (a command, duty, etc.); to disregard. Also, to neglect (to do something). \textit{Obsolete.}
				\item[3] To give up, renounce.
				\item[3.a] To give up, part with, surrender (esp. something dear or valued). Passing into sense 4.
				\item[3.b] To break off from, renounce (an employment, design, esp. an evil practice or sin; also, a belief, doctrine). 
				\item[4.a] To abandon, leave entirely, withdraw from; esp. to withdraw one's presence and help or companionship from; to desert.
				\item[4.b] Of things: To fail, disappoint the hopes of. \textit{Obsolete.}]
			\end{compactitem}

	% ----------------------
	\subsection{Futurity} % *FUTURITY
	% ----------------------
		\label{FUTURITY}
		
		% -----
		\subsubsection{1828 American Dictionary of the English Language}
		% -----
		
			\begin{compactitem}
				\item Future time; time to come.
				\item[1] Event to come.
				\item[2] The state of being yet to come, or to come hereafter.
			\end{compactitem}

% -----------------------------------------------------------
% -----------------------------------------------------------
% *G --------------------------------------------------------
% ----------------------------------------------------------- 
% -----------------------------------------------------------

	% ----------------------
	\subsection{Gad} % *GAD
	% ----------------------
		\label{GAD}
		
		% -----
		\subsubsection{NIEBC, ``Gad'' \#1}
		% -----
		
			``Gad was Jacob's seventh son and born to Leah's servant, Zilpah. He was therefore considered to be Leah's son (Gen. 30:11; 35:26). He was among those who went down to Egypt with his father and his many sons (Gen. 46:16). In Jacob's blessing he prophesied that Gad would be the subject of attack by raiders but would fight back (Gen. 49:19).
			
			``He became the leader of one of the twelve tribes of Israel named after him. The tribe went on to inherit land to the east of Jordan when finally the conquest of Canaan took place (Num. 32). At the Exodus from Egypt the fighting men in the tribe were numbered as 46,650 (Num. 1:25). The tribe of Gad is mentioned again in Ezek. 48 and Rev. 7:5 in texts that look forward to the last days and the fulfilment of the Lord's divine purposes for his kingdom.''
			
	% ----------------------
	\subsection{Gall} % *GALL
	% ----------------------
		\label{GALL}
		
		% -----
		\subsubsection{1828 American Dictionary of the English Language}
		% -----
		
			\begin{compactitem}
				\item[1] In the animal economy, the bile, a bitter, a yellowish green fluid, secreted in the glandular substance of the liver. It is glutinous or imperfectly fluid, like oil.
				\item[2] Any thing extremely bitter.
				\item[3] Rancor; malignity.
				\item[4] Anger; bitterness of mind.
			\end{compactitem}
			
		% -----
		\subsubsection{Oxford English Dictionary}
		% -----
		
			\begin{compactitem}
				\item[I.1.a] The secretion of the liver, bile.
				\item[I.1.b] \textit{figurative.} With reference to the bitterness of gall. \textbf{to dip one's pen in gall}, to write with virulence and rancour.
				\item[I.1.c] In Biblical phrases.
				\item[I.2.a] The gall bladder and its contents.
				\item[I.2.b] Short for `sickness of the gall', a disease in cattle. \textit{Obsolete.}
				\item[I.3.a] Bitterness of spirit, asperity, rancour (supposed to have its seat in the gall).
				\item[I.3.b] Spirit to resent injury or insult. \textit{Obsolete.}
				\item[I.3.c] Hence, \textbf{to break one's gall}: in early use, to break the spirit, cow, subdue.
					\begin{compactitem}
						\item a1500--1785
					\end{compactitem}
				\item[I.4] Assurance, impudence.
					\begin{compactitem}
						\item This is the sense of, ``You have the gall to tell me...''
					\end{compactitem}
				\item[II] In certain transferred uses.
				\item[II.5] Poison, venom. \textit{Obsolete.}
				\item[II.6] \textbf{gall of the earth}: a name given to the Lesser Centaury, from its bitterness.
				\item[II.7] The scum of melted glass.
			\end{compactitem}
			
	% ----------------------
	\subsection{Garnish} % *GARNISH
	% ----------------------
		\label{GARNISH}
		
		% -----
		\subsubsection{1828 American Dictionary of the English Language}
		% -----
		
			\begin{compactitem}
				\item[1] To adorn; to decorate with appendages; to set off.
				\item[2] To fit with fetters; a cant term.
				\item[3] To furnish; to supply; as a fort garnished with troops.
				\item[4] In law, to warn; to give notice. [See Garnishee.]
			\end{compactitem}
			
	% ----------------------
	\subsection{Graduation} % *GRADUATION
	% ----------------------
		\label{GRADUATION}
		
		% -----
		\subsubsection{See also}
		% -----
		
			\hyperref[DEGREE]{Degree}, 
		
		% -----
		\subsubsection{1828 American Dictionary of the English Language}
		% -----
		
			\begin{compactitem}
				\item[1] Improvement; exaltation of qualities.
				\item[2] The act of conferring or receiving academical degrees.
				\item[3] The act of marking with degrees.
				\item[4] The process of bringing a liquid to a certain consistence by evaporation.
			\end{compactitem}
		
		% -----
		\subsubsection{Oxford English Dictionary}
		% -----	
		
			\begin{compactitem}
				\item[1.a] The action or process of dividing into degrees or other proportionate divisions on a graduated scale.
				\item[1.b] \textit{pl.} Lines employed to indicate degrees of latitude and longitude, quantity, etc.
				\item[1.c] The manner in which something is graduated.
				\item[1.d] Position on a map as indicated by degrees. \textit{Obsolete.}
				\item[2.a] Arrangement in degrees or gradations; `regular progression by succession of degrees'
				\item[2.b] An elevation by degrees into a higher condition; also quasi-\textit{concr.} a step in the process, a degree.
				\item[3] \textit{Alch.}, \textit{Chem.}, etc.
				\item[3.a] The process of tempering the composition of a substance to a required degree; the process of refining an element, a metal. \textit{Obsolete.}
				\item[3.b] The process of concentrating by evaporation.
				\item[4] \textit{Gunnery.}
				\item[5] \textit{U. S. Railways.} Formerly used for `grading', `gradient'.
				\item[6] The action of receiving or conferring a university degree, or a certificate of qualification from some recognized authority. Also, the ceremony of conferring degrees.
			\end{compactitem}

% -----------------------------------------------------------
% -----------------------------------------------------------
% *H --------------------------------------------------------
% ----------------------------------------------------------- 
% -----------------------------------------------------------

	% ----------------------
	\subsection{Hereafter} % *HEREAFTER
	% ----------------------
		\label{HEREAFTER}
		
		From the 1828, the complete entry:
		
			\begin{itemize}
				\item In time to come; in some future time.
				\item[1] In a future state.
			\end{itemize}

	% ----------------------
	\subsection{Horeb (Mount)} % *HOREB
	% ----------------------
		\label{HOREB}
		
		See \hyperref[SINAI]{Sinai}.
		
	% ----------------------
	\subsection{Howbeit} % *HOWBEIT
	% ----------------------
		\label{HOWBEIT}
		
		From the 1828 dictionary for ``howbeit,'' the complete entry:
		
		\begin{compactitem}
			\item Be it as it may; nevertheless; notwithstanding; yet; but; however.
		\end{compactitem}

% -----------------------------------------------------------
% -----------------------------------------------------------
% *I --------------------------------------------------------
% ----------------------------------------------------------- 
% -----------------------------------------------------------

	% ----------------------
	\subsection{Immediate} % *IMMEDIATE
	% ----------------------
		\label{IMMEDIATE}
		
		% -----
		\subsubsection{1828 American Dictionary of the English Language}
		% -----
		
			\begin{compactitem}
				\item[1] Proximate; acting without a medium, or without the intervention of another cause or means; producing its effect by its own direct agency. An \textit{immediate} cause is that which is exerted directly in producing its effect, in opposition to a mediate cause, or one more remote.
				\item[2] Not acting by second causes; as the \textit{immediate} will of God.
				\item[3] Instant; present; without the intervention of time. We must have an \textit{immediate} supply of bread.
			\end{compactitem}

	% ----------------------
	\subsection{Interest} % *INTEREST
	% ----------------------
		\label{INTEREST}
		
		% -----
		\subsubsection{1828 American Dictionary of the English Language}
		% -----

			\begin{compactitem}
				\item Concern; advantage; good; as private \textit{interest}; public \textit{interest}
				\item[1] Influence over others. They had now lost their \textit{interest} at court.
				\item[2] Share; portion; part; participation in value. He has parted with his \textit{interest} in the stocks. He has an \textit{interest} in a manufactory of cotton goods.
				\item[3] Regard to private profit.
				\item[4] Premium paid for the use of money; the profit per cent derived from money lent, or property used by another person, or from debts remaining unpaid. Commercial states have a legal rate of \textit{interest} Debts on book bear an \textit{interest} after the expiration of the credit. Courts allow \textit{interest} in many cases where it is not stipulated. A higher rate of \textit{interest} than that which the law allows, is called usury.
				\item[5] Any surplus advantage.
			\end{compactitem}

% -----------------------------------------------------------
% -----------------------------------------------------------
% *J --------------------------------------------------------
% ----------------------------------------------------------- 
% -----------------------------------------------------------

	% ----------------------
	\subsection{Jethro} % *JETHRO
	% ----------------------
		\label{JETHRO}
	
		% -----
		\subsubsection{NIEBC, ``Jethro''}
		% -----
		
			``(Heb., `excellence'). Jethro was Moses' father-in-law. He was `a priest of Midian' and also known as Reuel (Exod. 3:1). He gave his daughter Zipporah to Moses to be his wife. Moses had helped Jethro's seven daughters water their flocks and so had become a close family friend. Having fled Egypt, Moses stayed with the family and helped them as a shepherd for many years until eventually he returned to Egypt to confront Pharaoh (Exod. 2:15--22; 4:18).

			``When eventually the Israelites escaped across the Red Sea into the wilderness, Jethro took Zipporah and they visited Moses in the desert (Exod. 18:1--27). There Jethro was encouraged by what he learned of all that the LORD had done for the Israelites and he clearly grew in the faith of Yahweh, saying to Moses: `Now I know that the LORD is greater than all other gods, for he did this to those who had treated Israel arrogantly' (Exod. 18:11).

			``It is of particular interest that this priest was obviously a priest to Yahweh. In a sacrifice of thanksgiving, at which Aaron was present, Jethro performed the sacrifice and Aaron and the other elders came and ate bread and had fellowship with him (v. 12).

			``During Jethro's stay in the camp of the Israelites he was worried about the amount of responsibility and work that Moses had. He was afraid Moses would be worn out. In a great display of wisdom, Jethro advocated a shared leadership in which Moses would fulfil the main task of coming before the LORD while other `capable men', who feared God and were trustworthy, would be selected to help judge the land for the people. They would decide simple legal cases between people and only the worst cases would be brought to Moses (vv. 17--26).

			``This wise and godly man recognized the great faithfulness of the LORD and helped God's people in learning how to govern themselves and in employing a plurality leaders.''

% -----------------------------------------------------------
% -----------------------------------------------------------
% *K --------------------------------------------------------
% ----------------------------------------------------------- 
% -----------------------------------------------------------

% -----------------------------------------------------------
% -----------------------------------------------------------
% *L --------------------------------------------------------
% ----------------------------------------------------------- 
% -----------------------------------------------------------

	% ----------------------
	\subsection{Luster} % *LUSTER
	% ----------------------
		\label{LUSTER}
		
		% -----
		\subsubsection{1828 American Dictionary of the English Language}
		% -----

			\begin{compactitem}
				\item[1] Brightness; splendor; gloss; as the \textit{luster} of the sun or stars; the \textit{luster} of silk.
				\item[2] The splendor of birth, of deeds or of fame; renown; distinction.
				\item[3] A sconce with lights; a branched candlestick of glass.
				\item[4] The space of five years. [Latin lustrum.]
			\end{compactitem}

% -----------------------------------------------------------
% -----------------------------------------------------------
% *M --------------------------------------------------------
% ----------------------------------------------------------- 
% -----------------------------------------------------------

	% ----------------------
	\subsection{Maintenance} % *MAINTENANCE
	% ----------------------
		\label{MAINTENANCE}
		
		% -----
		\subsubsection{1828 American Dictionary of the English Language}
		% -----
		
			\begin{compactitem}
				\item Sustenance; sustentation; support by means of supplies of food, clothing and other conveniences; as, his labor contributed little to the \textit{maintenance} of his family.
				\item[1] Means of support; that which supplies conveniences.
				\item[2] Support; protection; defense; vindication; as the \textit{maintenance} of right or just claims.
				\item[3] Continuance; security from failure or decline.
				\item[4] In law, an officious intermeddling in a suit in which the person has no interest, by assisting either party with money or means to prosecute or defend it. This is a punishable offense. But to assist a poor kinsman from compassion, is not \textit{maintenance}.
			\end{compactitem}

	% ----------------------
	\subsection{Meet} % *MEET
	% ----------------------
		\label{MEET}
		
		From the 1828 dictionary, the entry for ``meet'' as adjective:
		
			\begin{compactitem}
				\item Fit; suitable; proper; qualified; convenient; adapted, as to a use or purpose.
			\end{compactitem}

% -----------------------------------------------------------
% -----------------------------------------------------------
% *N --------------------------------------------------------
% ----------------------------------------------------------- 
% -----------------------------------------------------------

	% ----------------------
	\subsection{Notwithstanding} *NOTWITHSTANDING
	% ----------------------
		\label{NOTWITHSTANDING}
		
		% -----
		\subsubsection{1828 American Dictionary of the English Language}
		% -----
		
			\begin{compactitem}
				\item The participle of withstand, with not prefixed, and signifying not opposing; nevertheless. It retains in all cases its participial signification. For example, `I will surely rend the knigdom from thee, and will give it to thy servant; \textit{notwithstanding} in thy days I will not do it, for david thy father's sake.' 1 Kings 11:12. In this passage there is an ellipsis of that, after \textit{notwithstanding} That refers to the former part of the sentence, I will rend the kingdom from thee; \textit{notwithstanding} that (declaration or determination,) in thy days I will not do it. in this and in all cases, \textit{notwithstanding} either with or without that or this, constitutes the case absolute or independent.
				\item `It is a rainy day, but \textit{notwithstanding} that, the troops must be reviewed;' that is, the rainy day not opposing or preventing. That, in this case, is a substitute for the whole first clause of the sentence. It is to that clause what a relative is to an antecedent noun, and which may be used in the place of it; \textit{notwithstanding} which, that is, the rainy day.
				\item 'Christ enjoined on his followers not to publish the cures he wrought; but \textit{notwithstanding} his injunctions, they proclaimed them.' Here, \textit{notwithstanding} his injunctions, is the case independent or absolute; the injunctions of Christ not opposing or preventing.
				\item This word answers precisely to the latin \textit{non obstante}, and both are used with nouns or with substitutes for nouns, for sentences or for clauses of sentences. So in the Latin phrase, hoc non obstante, hoc may refer to a single word, to a sentence or to a series of sentences.
			\end{compactitem}

% -----------------------------------------------------------
% -----------------------------------------------------------
% *O --------------------------------------------------------
% ----------------------------------------------------------- 
% -----------------------------------------------------------

% -----------------------------------------------------------
% -----------------------------------------------------------
% *P --------------------------------------------------------
% ----------------------------------------------------------- 
% -----------------------------------------------------------

	% ----------------------
	\subsection{Persuasion} % *PERSUASION
	% ----------------------
		\label{PERSUASION}
		
		% -----
		\subsubsection{1828 American Dictionary of the English Language}
		% -----
		
			\begin{compactitem}
				\item[1] The act of persuading; the act of influencing the mind by arguments or reasons offered, or by any thing that moves the mind or passions, or inclines the will to a determination.
				\item[2] The state of being persuaded or convinced; settled opinion or conviction proceeding from arguments and reasons offered by others, or suggested by one's own reflections.
				\item[3] A creed or belief; or a sect or party adhering to a creed or system of opinions; as men of the same persuasion; all persuasions concur in the measure.
			\end{compactitem}

	% ----------------------
	\subsection{Precept} % *PRECEPT
	% ----------------------
		\label{PRECEPT}
		
		% -----
		\subsubsection{1828 American Dictionary of the English Language}
		% -----
		
			\begin{compactitem}
				\item[1] In a general sense, any commandment or order intended as an authoritative rule of action; but applied particularly to commands respecting moral conduct. The ten commandments are so many precepts for the regulation of our moral conduct.
				\item[2] In law, a command or mandate in writing.
			\end{compactitem}
			
	% ----------------------
	\subsection{Preparation} % *PREPARATION
	% ----------------------
		\label{PREPARATION}
		
		% -----
		\subsubsection{1828 American Dictionary of the English Language}
		% -----
		
			\begin{compactitem}
				\item[1] The act or operation of preparing or fitting for a particular purpose, use, service or condition.
				\item[2] Previous measures of adaptation.
				\item[3] Ceremonious introduction. [Unusual.]
				\item[4] That which is prepared, made or compounded for a particular purpose.
				\item[5] The state of being prepared or in readiness.
				\item[6] Accomplishment; qualification.
				\item[7] In pharmacy, any medicinal substance fitted for the use of the patient.
				\item[8] In anatomy, the parts of animal bodies prepared and preserved for anatomical uses.
			\end{compactitem}
			
	% ----------------------
	\subsection{Principal} % *PRINCIPAL
	% ----------------------
		\label{PRINCIPAL}
		
		% -----
		\subsubsection{1828 American Dictionary of the English Language}
		% -----	
		
			% --
			\paragraph{Noun}
			% --
		
				\begin{compactitem}
					\item A chief or head; one who takes the lead; as the \textit{principal} of a faction, an insurrection or mutiny.
					\item[1] The president, governor, or chief in authority. We apply the word to the chief instructor of an academy or seminary of learning.
					\item[2] In law, the actor or absolute perpetrator of a crime, or an abettor. A \textit{principal} in the first degree, is the absolute perpetrator of the crime; a \textit{principal} in the second degree, is one who is present, aiding and abetting the fact to be done; distinguished from an accessory. In treason, all persons concerned are \textit{principal}s.
					\item[3] In commerce, a capital sum lent on interest, due as a debt or used as a fund; so called in distinction from interest or profits.
					\item[4] One primarily engaged; a chief party; in distinction from an auxiliary.
				\end{compactitem}
			
	% ----------------------
	\subsection{Propitiation} % *PROPITIATION
	% ----------------------
		\label{PROPITIATION}
		
		% -----
		\subsubsection{1828 American Dictionary of the English Language}
		% -----
		
			% --
			\paragraph{Propitiation}
			% --
		
				\begin{compactitem}
					\item[1] The act of appeasing wrath and conciliating the favor of an offended person; the act of making propitious.
					\item[2] In theology, the atonement or atoning sacrifice offered to God to assuage his wrath and render him propitious to sinners. Christ is the \textit{propitiation} for the sins of men. Romans 3:25. 1 John 2:2.
				\end{compactitem}
				
			% --
			\paragraph{Propitious}
			% --
			
				\begin{compactitem}
					\item[1] Disposed to be gracious or merciful; ready to forgive sins and bestow blessings; applied to God.
					\item[2] Favorable; as a \textit{propitious} season.
				\end{compactitem}
	
% -----------------------------------------------------------
% -----------------------------------------------------------
% *Q --------------------------------------------------------
% -----------------------------------------------------------
% -----------------------------------------------------------
		
% -----------------------------------------------------------
% -----------------------------------------------------------
% *R --------------------------------------------------------
% ----------------------------------------------------------- 
% -----------------------------------------------------------

	% ----------------------
	\subsection{Redound} % *REDOUND
	% ----------------------
		\label{REDOUND}
	
		% -----
		\subsubsection{1828 American Dictionary of the English Language}
		% -----

			\begin{compactitem}
				\item[1] To be sent, rolled or driven back.
				\item[2] To conduce in the consequence; to contribute; to result.
				\item[3] To proceed in the consequence or effect; to result.
			\end{compactitem}

	% ----------------------
	\subsection{Rundle} % *RUNDLE
	% ----------------------
		\label{RUNDLE}
	
		% -----
		\subsubsection{1828 American Dictionary of the English Language}
		% -----

			\begin{compactitem}
				\item[1] A round; a step of a ladder.
				\item[2] Something put round an axis.
			\end{compactitem}

% -----------------------------------------------------------
% -----------------------------------------------------------
% *S --------------------------------------------------------
% ----------------------------------------------------------- 
% -----------------------------------------------------------

	% ----------------------
	\subsection{Sabaoth (Lord of)} % *SABAOTH
	% ----------------------
		\label{SABAOTH}
		
		This word occurs twice in the NT (\hyperref[ROMANS-9-29]{Romans 9:29}, \hyperref[JAMES-5-4]{James 5:4} and four times in the DC (\hyperref[DC87-7]{DC 87:7}, \hyperref[DC88-2]{88:2}, \hyperref[DC95-7]{95:7}, \hyperref[DC98-2]{98:2}). In the NT it is \gr{σαβαώθ}, which is a transliteration of \he{.sAbA'}, or more correctly the plural \he{.s:bA'Ot}. In the OT, the KJV translates \he{.sAbA'} and its various forms as ``host'' (the majority of occurrences), but also ``war,'' ``army,'' `battle,'' and a few other terms.
		
		Note \hyperref[DC95-7]{DC 95:7} especially: ``And for this cause I gave unto you a commandment that you should call your solemn assembly, that your fastings and your mourning might come up into the ears of the Lord of Sabaoth, which is by interpretation, the creator of the first day, the beginning and the end.'' I have no idea what to say about this; is it possible that JS saw it as something to do with ``Sabbath,'' and so it had to do with days?
		
		% -----
		\subsubsection{\texorpdfstring{\he{.sAbA'}}{}} % צָבָא
		% -----
			\label{TSAVA} % whatever is in step bible without periods
		
			\begin{compactdesc}
				% \item[HALOT , PDF ] \textbf{} % Use the 1994 PDF
					% \begin{compactitem}
						% \item[1]
					% \end{compactitem}
				% \item[BDB , PDF ] \textbf{}
					% \begin{compactitem}
						% \item[1]
					% \end{compactitem}
				% \item[DCH PDF ] \textbf{} % don't include the actual page number, they repeat from volume to volume
					% \begin{compactitem}
						% \item[1]
					% \end{compactitem}
				\item[TWOT 1865b, 750]	\phantom{.}
					\begin{compactitem}
						\item ``Translated as host(s) \he{.sAbA'} means army(ies). It can refer to any arrayed army (\hyperref[JUDGES-4-2]{Jud 4:2}), the inhabitants of heaven (\hyperref[1KINGS-22-19]{I Kgs 22:19}), or the celestial bodies (\hyperref[DEUTERONOMY-4-19]{Deut 4:19}).''
						\item ``The heavenly bodies, including the sun and the moon, are called the host of heaven (\hyperref[GENESIS-2-1]{Gen 2:1}). When referring to them the word is always singular. God created this host by his breath (\hyperref[PSALMS-33-6]{Ps 33:6}), and he preserves their existence (\hyperref[ISAIAH-40-26]{Isa 40:26}). They thus serve and worship him in complete obedience (\hyperref[NEHEMIAH-9-6]{Neh 9:6}; \hyperref[ISAIAH-45-12]{Isa 45:12}). The host are identified as his ministers that do his will (\hyperref[PSALMS-103-21]{Ps 103:21}). Israel is exhorted never to worship them (\hyperref[DEUTERONOMY-4-19]{Deut 4: 19}).''
						\item ``This divine name [Yahweh and \he{'E:lohiyM}] appears for the first time in I Sam 1:3. Its origin appears to have been at the close of the period of the judges and in the vicinity ofthe sanctuary Shiloh, where the ark of the covenant was housed. The ark itself symbolized Yahweh's rulership; for he is declared to be enthroned between the cherubim (I Sam 4:4; cf. Ps 99:1). This name certainly contains the affirmation that Yahweh is the true head of lsrael's armies. The idea that more than Israel's armies is encompassed in this title is clear from David's statement, `Yahweh of hosts, the God of the armies of Israel' (I Sam 17:45). Rather it affirms his universal rulership that encompasses every force or army, heavenly, cosmic and earthly. Now that Israel was emerging as a nation with international relationships, the language which exposed the theology of its God needed to keep pace. It was important to affirm that Yahweh was not merely one warrior god among the leading warrior gods of the nations, but that he was the Supreme God. Particularly for Israel, located on the landbridge between three major continents which was constantly crossed by the armies of the great world powers, it became essential to emphasize that Yahweh was King even of the armies of these mighty empires. As a result he was sufficient to lead Israel to overcome any crisis brought on by those armies. The prophets, during the kingdom period, faced a further problem, namely that God used these mighty armies to punish Israel in their rebellion against God. Hence it was essential to point out that Yahweh indeed was the king of those nations and that he would judge them. Conversely, if Israel would not return to God, then Yahweh could employ those armies against her and reduce her to captivity (cf. Isa 10:5--34). Further for the prophets during the time of the Babylonian crises when the people were attracted to astral worship, this name conveyed clearly to the people that it was foolish to worship these stars which were merely obedient creatures to Yahweh (cf. Isa 47).
						\item ``When captured in all of its thrust the name \he{yhwh .s:bA'Ot} is a most exalted title. It is definitely associated with Yahweh's kingship as Isa 6:5 and Ps 84:3 [H 4] show. On a festive day before a triumphal procession entered the temple court, the chorus sang: `Lift up your heads, O gates! and be lifted up, O ancient doors! that the King of glory may come in. Who is this King of glory? The Lord of hosts, he is the King of glory!' (Ps 24:9f.). The text here clearly shows that Yahweh of hosts conveys the concept of glorious king. Yahweh is King of the world (cf. Zech 14:16) and over all the kingdoms of the earth (Isa 37:16). This God is the source of all (Jer 10:16). His rulership necessitates a time when he will visibly display that lordship. In the last days the nations will wage war against Mount Zion. They will think they have won, only to come to the realization that their imagination has outdistanced their accomplishment. At that time Yahweh will lead a great army into battle aided by the forces of nature. Every opponent in heaven and on earth will be subjugated (Isa 13:4; 24:21ff.; 29:5--8; 31:4f.; 34: 1--12). Then Yahweh will visibly manifest a universal reign from Mt. Zion. All kings and nations will acknowledge that the Yahweh of hosts is the king of glory. His total authority and Iordship will be etemally established. 
						\item ``Although the title has military overtones, it points directly to Yahweh's rulership over the entire universe. He continually rules, but at times he directly intervenes to secure his own victory and insure the direction of history for the salvation of his people. In Amos 4:13 it is associated with his creating the mountains and wind and his ability to control nature (cf. Amos 5:8f.; 9:5f.). He is master over every force; he alone secures peace. To him prayer may be addressed (e.g. Ps 80:19 [H 20]). Special attention is given to the majestic splendor of Yahweh's rule in this title (cf. Ps 84:1 [H 2]; Isa 28:5f.).''
					\end{compactitem}
				\item[TDOT 12:215, PDF 236] \phantom{.}
					% \begin{compactitem}
						% \item 
					% \end{compactitem}
			\end{compactdesc}
			
	% ----------------------
	\subsection{Shrink} % *SHRINK
	% ----------------------
		\label{SHRINK}
	
		% -----
		\subsubsection{1828 English Dictionary}
		% -----
		
			\begin{compactitem}
				\item[1] To contract spontaneously; to draw or be drawn into less length, breadth or compass by an inherent power; as, woolen cloth \textit{shrinks} in hot water; a flaxen of hempen line \textit{shrinks} in a humid atmosphere. Many substances \textit{shrink} by drying.
				\item[2] To shrivel; to become wrinkled by contraction; as the skin.
				\item[3] To withdraw or retire, as from danger; to decline action from fear. A brave man never \textit{shrinks} from danger; a good man does not \textit{shrink} from duty.
				\item[4] To recoil, as in fear, horror or distress. My mind \textit{shrinks} from the recital of our woes.
				\item[5] To express fear, horror or pain by shrugging or contracting the body.
			\end{compactitem}
			
	% ----------------------
	\subsection{Similitude} % *SIMILITUDE
	% ----------------------
		\label{SIMILITUDE}
		
		% -----
		\subsubsection{1828 English Dictionary}
		% -----
		
			\begin{compactitem}
				\item[1] Likeness; resemblance; likeness in nature, qualities of appearance; as \textit{similitude} of substance. Let us make man in our image, man in our \textit{similitude}. Fate some future bard shall join in sad \textit{similitude} of griefs to mine.
				\item[2] Comparison; simile. Tasso, in his \textit{similitude} never departed from the woods.
			\end{compactitem}

	% ----------------------
	\subsection{Sinai (Mount)} % *SINAI
	% ----------------------
		\label{SINAI}
		
		There's a scholarly consensus that Mount Sinai and Mount Horeb are the same.
		
	% ----------------------
	\subsection{Sociality} % *SOCIALITY
	% ----------------------
		\label{SOCIALITY}
	
		% -----
		\subsubsection{Oxford English Dictionary}
		% -----
		
			\begin{compactitem}
				\item[1.a] The state or quality of being sociable; (the enjoyment of) friendly social interaction; sociability.
				\item[1.b] Friendship, companionship, or social interaction with a person or (occasionally) in a certain state; an instance of this. Also in extended use.
				\item[1.c] Formal or conventional social interaction. Contrasted with sociability. \textit{Obsolete. rare.}
				\item[2.a] The disposition or tendency to live in a society or to associate with others.
				\item[2.b] Social organization or behaviour in animals; the tendency to live, or the fact of living, in a community of individuals of the same species which cooperate with one another to their mutual or collective benefit.
				\item[3] In \textit{plural}. Social pleasantries, events, or entertainments.
			\end{compactitem}
			
	% ----------------------
	\subsection{Somewhat} % *SOMEWHAT
	% ----------------------
		\label{SOMEWHAT}
	
		% -----
		\subsubsection{1828 American Dictionary of the English Language}
		% -----
		
			% --
			\paragraph{Noun}
			% --
			
				\begin{compactitem}
					\item[1] Something, though uncertain what.
					\item[2] More or less; a certain quantity or degree, indeterminate.
					\item[3] A part, greater or less.
				\end{compactitem}
		
			% --
			\paragraph{Adverb}
			% --
		
				\begin{compactitem}
					\item In some degree or quantity.
				\end{compactitem}
				
	% ----------------------
	\subsection{Straight} % *STRAIGHT
	% ----------------------
		\label{STRAIGHT}
	
		% -----
		\subsubsection{1828 American Dictionary of the English Language}
		% -----
		
			\begin{compactitem}
				\item Latin, formed from the root of reach, stretch, right. It is customary to write \textit{straight} for direct or right, and \textit{strait}, for narrow, but this is a practice wholly arbitrary, both being the same word. Strait we use in the sense in which it is used in the south of Europe. Both sense proceed from stretching, straining.
				\item[1] Right, in a mathematical sense; direct; passing from one point to another by the nearest course; not deviating or crooked; as a \textit{straight} line; a \textit{straight} course; a \textit{straight} piece of timber.
				\item[2] Narrow; close; tight; as a \textit{straight} garment. [See strait, as it is generally written.]
				\item[3] Upright; according with justice and rectitude; not deviating from truth or fairness.
			\end{compactitem}
				
	% ----------------------
	\subsection{Strait} % *STRAIT
	% ----------------------
		\label{STRAIT}
	
		% -----
		\subsubsection{1828 American Dictionary of the English Language}
		% -----
		
			\begin{compactitem}
				\item[1] Narrow; close; not broad.
				\item[2] Close; intimate; as a strait degree of favor.
				\item[3] Strict; rigorous.
				\item[4] Difficult; distressful.
				\item[5] Straight; not crooked.
			\end{compactitem}

% -----------------------------------------------------------
% -----------------------------------------------------------
% *T --------------------------------------------------------
% ----------------------------------------------------------- 
% -----------------------------------------------------------

% -----------------------------------------------------------
% -----------------------------------------------------------
% *U --------------------------------------------------------
% ----------------------------------------------------------- 
% -----------------------------------------------------------

	% ----------------------
	\subsection{Unspeakable} % *UNSPEAKABLE
	% ----------------------
		\label{UNSPEAKABLE}
		
		% -----
		\subsubsection{1828 American Dictionary of the English Language}
		% -----

			\begin{compactitem}
				\item That cannot be uttered; that cannot be expressed; unutterable; as \textit{unspeakable} grief or rage.
			\end{compactitem}

% -----------------------------------------------------------
% -----------------------------------------------------------
% *V --------------------------------------------------------
% ----------------------------------------------------------- 
% -----------------------------------------------------------

	% ----------------------
	\subsection{Valiant} % *VALIANT
	% ----------------------
		\label{VALIANT}
		
		% -----
		\subsubsection{1828 American Dictionary of the English Language}
		% -----
		
			\begin{compactitem}
				\item[1] Primarily, strong; vigorous in body; as a \textit{valiant} fencer.
				\item[2] Brave; courageous; intrepid in danger; heroic; as a \textit{valiant} soldier.
				\item[3] Performed with valor; bravely conducted; heroic; as a \textit{valiant} action or achievement; a \textit{valiant} combat.
			\end{compactitem}

% -----------------------------------------------------------
% -----------------------------------------------------------
% *W --------------------------------------------------------
% ----------------------------------------------------------- 
% -----------------------------------------------------------

	% ----------------------
	\subsection{Wherefore} % *WHEREFORE
	% ----------------------
		\label{WHEREFORE}
		
		% -----
		\subsubsection{1828 American Dictionary of the English Language}
		% -----

			\begin{compactitem}
				\item[1] for which reason
				\item[2] why; for what reason
			\end{compactitem}

		% -----
		\subsubsection{Oxford English Dictionary} % *OED
		% -----
			
			When it isn't used as an interrogative, the most important meanings for the scriptures are II.4 and II.5.a, which are actually pretty similar.

			\begin{compactitem}
				\item[I] Interrogative uses
				\item[I.1] For what? \textit{esp.} for what purpose or end?
				\item[I.2] For what cause or reason? on what account? why?
				\item[II] Relative uses
				\item[II.3] For which
				\item[II.4] On account of or because of which; in consequence or as a result of which. Chiefly with noun (esp. \textit{reason} or \textit{cause}) as antecedent. \textit{archaic}.
				\item[II.5.a] Introducing a clause expressing a consequence or inference from what has just been stated: On which account; for which reason; which being the case; and therefore.
			\end{compactitem}
			
	% ----------------------
	\subsection{Whoso} % *WHOSO
	% ----------------------
		\label{WHOSO}
		
		% -----
		\subsubsection{1828 American Dictionary of the English Language}
		% -----

			\begin{compactitem}
				\item Any person whatever.
			\end{compactitem}

	% ----------------------
	\subsection{Wing} % *WING
	% ----------------------
		\label{WING}
		
		% -----
		\subsubsection{1828 American Dictionary of the English Language}
		% -----
		
			\begin{compactitem}
				\item[1] The limb of a fowl by which it flies. In a few species of fowls, the wings do not enable them to fly; as is the case with the dodo, ostrich, great auk, and penguin; but in the two former, the wings assist the fowls in running.
				\item[2] The limb of an insect by which it flies.
				\item[3] In botany, the side petal of a papilionaceous corol; also, an appendage of seeds, by means of which they are wafted in the air and scattered; also, any membranous or leafy dilatation of a footstalk, or of the angles of a stem, branch or flower stalk, or of a calyx.
				\item[4] Flight; passage by the wind; as, to be on the wind; to take wing.
				\item[5] Means of flying; acceleration. Fear adds wings to flight.
				\item[6] Motive or incitement of flight.
				\item[7] The flank or extreme body or part of an army.
				\item[8] Any side-piece.
				\item[9] In gardening, a side-shoot.
				\item[10] In architecture, a side-building, less than the main edifice.
				\item[11] In fortification, the longer sides of hornworks, crown-works, etc.
				\item[12] In a fleet, the ships on the extremities, when ranged in a line, or when forming the two sides of a triangle.
				\item[13] In a ship, the wings are those parts of the hold and orlop deck, which are nearest the sides.
				\item[14] In Scripture, protection; generally in the plural. Psalms 63:7. Exodus 19:4.
			\end{compactitem}

% -----------------------------------------------------------
% -----------------------------------------------------------
% *X --------------------------------------------------------
% ----------------------------------------------------------- 
% -----------------------------------------------------------

% -----------------------------------------------------------
% -----------------------------------------------------------
% *Y --------------------------------------------------------
% ----------------------------------------------------------- 
% -----------------------------------------------------------

% -----------------------------------------------------------
% -----------------------------------------------------------
% *Z --------------------------------------------------------
% ----------------------------------------------------------- 
% -----------------------------------------------------------